% !TeX root = ../tesi.tex

\capitolononnumerato{Acronimi e abbreviazioni}

Elenco in ordine alfabetico gli acronimi e le abbreviazioni introdotti nei capitoli 1, 2 e 3. Quando
presente, il collegamento rimanda all’approfondimento nel glossario.

\rimandoglossario{adr}{\textbf{ADR}}: \emph{Architecture Decision Record}.

\rimandoglossario{api}{\textbf{API}}: \emph{Application Programming Interface}.

\textbf{AWS}: Amazon Web Services.

\rimandoglossario{b2b}{\textbf{B2B}}: \emph{Business-to-Business}.

\rimandoglossario{b2c}{\textbf{B2C}}: \emph{Business-to-Consumer}.

\textbf{CLI}: \emph{Command Line Interface}.

\textbf{CMS}: \emph{Content Management System}.

\rimandoglossario{crm}{\textbf{CRM}}: \emph{Customer Relationship Management}.

\textbf{CSS}: \emph{Cascading Style Sheets}.

\rimandoglossario{erp}{\textbf{ERP}}: \emph{Enterprise Resource Planning}.

\textbf{FAQ}: \emph{Frequently Asked Questions}.

\textbf{HTML}: \emph{HyperText Markup Language}.

\textbf{HTTP}: \emph{HyperText Transfer Protocol}.

\textbf{IA}: intelligenza artificiale.

\rimandoglossario{ipaas}{\textbf{iPaaS}}: \emph{Integration Platform as a Service}.

\textbf{JSON}: \emph{JavaScript Object Notation}.

\textbf{JWT}: \emph{JSON Web Token}.

\rimandoglossario{llm}{\textbf{LLM}}: \emph{Large Language Model}.

\rimandoglossario{mcp}{\textbf{MCP}}: \emph{Model Context Protocol}.

\rimandoglossario{sal}{\textbf{SAL}}: stato di avanzamento dei lavori.

\rimandoglossario{semver}{\textbf{SemVer}}: \emph{Semantic Versioning}.

\rimandoglossario{sse}{\textbf{SSE}}: \emph{Server-Sent Events}.

\rimandoglossario{ucp}{\textbf{UCP}}: \emph{Universal Commerce Protocol}.

\textbf{UI}: \emph{User Interface}.

\textbf{UML}: \emph{Unified Modeling Language}.

\textbf{UX}: \emph{User Experience}.
